\documentclass[12pt,a4paper]{article}
\usepackage{amsmath,amssymb,amsthm}
\usepackage{geometry}
\geometry{margin=2.5cm}
\usepackage{hyperref}
\hypersetup{colorlinks=true,linkcolor=blue,citecolor=blue,urlcolor=blue}
\usepackage{booktabs}
\usepackage{microtype}
\usepackage{lmodern}
\usepackage[T1]{fontenc}

% Theorem environments
\newtheorem{theorem}{Theorem}[section]
\newtheorem{proposition}[theorem]{Proposition}
\newtheorem{lemma}[theorem]{Lemma}
\newtheorem{corollary}[theorem]{Corollary}
\newtheorem{definition}[theorem]{Definition}
\newtheorem*{remark}{Remark}
\newtheorem*{observation}{Observation}
\newtheorem*{hypothesis}{Hypothesis}

\title{\textbf{Geometric Origin of Lorentzian Spacetime\\via Gnomonic Projection:\\A Unified Description through\\the Discretization Constant $C$}}

\author{Noriaki Kihara\\
\small WF System Co., Ltd.\\
\small (Osaka University, School of Engineering Science, Alumni)\\
\small \href{https://orcid.org/0009-0004-6753-4020}{ORCID: 0009-0004-6753-4020}\\
\small \href{https://doi.org/10.5281/zenodo.19434932}{DOI: 10.5281/zenodo.19434932}}

\date{April 2026\\[5pt]\small Research Note (Geometric Investigation)}

\begin{document}

\maketitle

\begin{abstract}
In the previous paper~\cite{Paper1}, we mapped 4-dimensional Euclidean space onto the hypersphere $S^4(R)$ via gnomonic projection and derived the geometric identity $G_{\mu\nu} + \Lambda g_{\mu\nu} = 0$ with $\Lambda = 3/R^2$. However, the following questions remained open: (1) Why does time carry a metric signature different from space? (2) Can the mapping be defined on discrete coordinates? (3) How many free parameters are required?

This paper provides geometric answers. First, motivated by the unobservability of the $R$-direction, we axiomatically introduce the Lorentzian signature (Axiom~3.1). Second, we show that quadratic metrics are a number-theoretic requirement for discretization (\S6.1) and that a single coupling constant $C$ (dimension: length) suffices (\S6.2). We prove that global regularity is guaranteed under Euclidean signature, while a cosmological horizon naturally emerges under Lorentzian signature (Theorem~6.1). The number-theoretic anomaly at $N=3$ spatial dimensions and the necessity of a fifth dimension are established via Fermat's two-square theorem.
\end{abstract}

\tableofcontents

%========================================
\section{Introduction}
%========================================

\subsection{Results of the Previous Paper}

The previous paper~\cite{Paper1} considered the central projection (gnomonic projection) from $n$-dimensional Euclidean space onto the $(n+1)$-dimensional hypersphere $S^n(R)$:
\begin{equation}
\Phi: \Pi_R \to S^n(R), \quad Y^\mu = \frac{Rx^\mu}{l}, \quad Y^{n+1} = \frac{R^2}{l}, \quad l = \sqrt{R^2 + \sum_\mu (x^\mu)^2} \tag{1.1}
\end{equation}
For $n = 4$, the pullback metric was derived:
\begin{equation}
g_{\mu\nu} = \frac{R^2}{l^2}\left(\delta_{\mu\nu} - \frac{x_\mu x_\nu}{l^2}\right) \tag{1.2}
\end{equation}
and $S^4(R)$ satisfies:
\begin{equation}
G_{\mu\nu} + \Lambda g_{\mu\nu} = 0, \qquad \Lambda = \frac{3}{R^2} \tag{1.3}
\end{equation}
Appendix~D of~\cite{Paper1} introduced signature parameters $\epsilon_\mu \in \{+1, -1\}$:
\begin{equation}
g_{\mu\nu} = \frac{R^2}{l^2}\left(\epsilon_\mu \delta_{\mu\nu} - \frac{\epsilon_\mu \epsilon_\nu x_\mu x_\nu}{l^2}\right), \qquad l^2 = R^2 + \sum_\mu \epsilon_\mu (x^\mu)^2 \tag{1.4}
\end{equation}

\subsection{Open Questions}

The previous paper presented $\epsilon_4 = -1$ (Lorentzian signature) as an \emph{option}, but did not explain why it is physically realized. This paper provides a geometric motivation, systematically develops the consequences, and clarifies the motivation for discrete coordinates.

\subsection{Organization}

\S2: Einstein tensor in general dimensions. \S3: Lorentzian signature. \S4: Lorentz group. \S5: Cognitive limits. \S6: Discrete coordinates and $C$. \S7: Singularity avoidance. \S8: Consistency checks. \S9: Information preservation. \S10: Outlook.

%========================================
\section{Einstein Tensor in General Dimensions}
%========================================

\begin{proposition}[Einstein tensor in general dimensions~{\rm\cite{doCarmo1992,Wald1984}}]\label{prop:general}
On the $n$-dimensional sphere $S^n(R)$:
\begin{equation}
R_{\mu\nu\rho\sigma} = \frac{1}{R^2}(g_{\mu\rho}g_{\nu\sigma} - g_{\mu\sigma}g_{\nu\rho}), \qquad R_{\mu\nu} = \frac{n-1}{R^2}g_{\mu\nu} \tag{2.1}
\end{equation}
\begin{equation}
R_{\text{scalar}} = \frac{n(n-1)}{R^2}, \qquad G_{\mu\nu} = -\frac{(n-1)(n-2)}{2R^2}g_{\mu\nu} \tag{2.2}
\end{equation}
\begin{equation}
\Lambda_n = \frac{(n-1)(n-2)}{2R^2} \tag{2.3}
\end{equation}
\end{proposition}

\begin{center}
\begin{tabular}{ccl}
\toprule
Dimension $n$ & $\Lambda_n$ & Status \\
\midrule
1 & 0 & Trivial \\
2 & 0 & Identically zero \\
3 & $1/R^2$ & Non-trivial (minimum) \\
\textbf{4} & $\mathbf{3/R^2}$ & \textbf{Main result of~\cite{Paper1}} \\
\bottomrule
\end{tabular}
\end{center}

\begin{remark}
Non-trivial GR requires $n \geq 3$; $n = 4$ is physically distinguished.
\end{remark}

%========================================
\section{Geometric Motivation and Axiomatic Setting for the Lorentzian Signature}
%========================================

\subsection{Structural Singularity of the $R$-Direction}

The 5-dimensional image of the gnomonic mapping~(1.1):
\[
\underbrace{Y^1, Y^2, Y^3, Y^4}_{\text{4D coordinates}} = \frac{Rx^\mu}{l}, \qquad \underbrace{Y^5}_{\text{$R$-direction}} = \frac{R^2}{l}
\]

\begin{proposition}[Asymmetry of $R$-direction]\label{prop:asym}
$Y^5$ is qualitatively different from $Y^1, \ldots, Y^4$: the 4D coordinates are intrinsic to $S^4(R)$, while $Y^5$ is the normal direction in the embedding space and is in principle unmeasurable by an observer confined to the sphere.
\end{proposition}

\subsection{Imaginary Rotation of the Unobservable Dimension}

\begin{definition}[Axiom~3.1: Axiomatic setting of the Lorentzian signature]\label{ax:lor}
We adopt two axioms:

\textbf{Axiom A} (Cognitive constraint): The intrinsic metric $g_{\mu\nu}$ of $S^n(R)$ (eq.~(1.2)) does not contain the embedding coordinate $Y^5$ (Gauss's Theorema Egregium). An observer using only the intrinsic metric cannot directly observe the $Y^5$-direction.

\textbf{Axiom B} (Imaginary rotation of unobservable dimension): The $R$-direction, intrinsically unobservable by Axiom~A, is represented as an effective imaginary axis.
\end{definition}

Axiom~A follows from Gauss's Theorema Egregium and admits no room for assumption. Axiom~B is a physical assumption.

Under Axioms A and B:

\begin{enumerate}
\item[(i)] $x^1, x^2, x^3$ are identified as spatial, $x^4 = ct$ as temporal.
\item[(ii)] $w_{\text{eff}} = i/R$, so $w_{\text{eff}}^2 = -1/R^2$.
\item[(iii)] The sign of $x^4 = ct$ flips: $\epsilon_4 = +1 \to -1$.
\end{enumerate}

Result:
\begin{equation}
l^2 = R^2 + x^2 + y^2 + z^2 - c^2t^2 \tag{3.3}
\end{equation}
\textbf{The Lorentzian signature $(-,+,+,+)$ is obtained under Axioms~A and B.}

\begin{remark}[Status of this claim]
Axiom~B could be replaced by other interpretations (compactification, gauge freedom, etc.). This paper explores consequences of Axiom~B, justified by consistency with observation.
\end{remark}

\subsection{Relationship to Conventional Wick Rotation}

This paper reverses the causal direction: the Lorentzian world is the ``cognitive reality'' for internal observers, while the Euclidean structure is the underlying geometric entity.

%========================================
\section{Geometric Origin of the Lorentz Group}
%========================================

\begin{proposition}[Chain of symmetry groups~{\rm\cite{Wald1984,Guo2004}}]\label{prop:chain}
\begin{equation}
SO(5) \xrightarrow[\text{Axiom~3.1}]{\epsilon_4: +1 \to -1} SO(4,1) \xrightarrow[R \to \infty]{\text{flat limit}} SO(3,1) \tag{4.1}
\end{equation}
where $SO(5)$ is the isometry group of Euclidean $S^4(R)$, $SO(4,1)$ is the de Sitter group, and $SO(3,1)$ is the Lorentz group.
\end{proposition}

Lorentz transformations are geometrically derived as the flat limit of de Sitter symmetry. The constancy of the speed of light is a consequence; the more fundamental principle is the symmetry of the gnomonic mapping.

%========================================
\section{Cognitive Limits of Internal Observers}
%========================================

\begin{theorem}[Interior angle upper bound]\label{thm:angle}
For a geodesic triangle on $S^n(R)$:
\begin{equation}
\alpha + \beta + \gamma = \pi + \frac{A}{R^2} \tag{5.1}
\end{equation}
At maximum area $A_{\max} = 2\pi R^2$:
\begin{equation}
(\alpha + \beta + \gamma)_{\max} = 3\pi = 540^\circ \tag{5.2}
\end{equation}
This upper bound is a universal constant \textbf{independent of $R$}.
\end{theorem}

\begin{proof}
By the Gauss--Bonnet theorem~\cite{doCarmo1992}, $\alpha + \beta + \gamma - \pi = A/R^2$. Maximum $A = 2\pi R^2$ gives $3\pi$.
\end{proof}

\begin{remark}
$R$ itself is measurable via the angular excess $\Delta = A/R^2$ for known $A$.
\end{remark}

\begin{proposition}[Conformal factor]
The gnomonic mapping involves the position-dependent conformal factor $R/l$, with $R/l = 1$ at the origin and $R/l \to 0$ as $r \to \infty$.
\end{proposition}

%========================================
\section{Discrete Coordinates and the Coupling Constant $C$}
%========================================

\subsection{Number-Theoretic Motivation}

\begin{observation}[Number-theoretic motivation for quadratic forms~{\rm\cite{Hurwitz1898}}]
On the integer lattice $\mathbb{Z}^n$, a multiplicative norm ($N(zw) = N(z)N(w)$) taking integer values requires the Gaussian integer norm $N(a+bi) = a^2 + b^2$ as minimal structure. By Hurwitz's theorem~(1898), composition algebras with multiplicative norms exist only in dimensions $1, 2, 4, 8$.
\end{observation}

\begin{remark}[Circumstantial evidence for discreteness]
Quadratic metrics are consistent with continuous spacetime but \emph{required} for discretization:
\[
\text{Discreteness} \Longrightarrow \text{Quadratic metrics required}
\]
\[
\text{Quadratic metrics} \not\Longrightarrow \text{Discreteness (but serve as evidence)}
\]
\end{remark}

\subsection{Introduction of $C$}

\begin{definition}[Coupling constant $C$]\label{def:C}
A constant $C > 0$ with dimension of length defines $w = C/R$. Under the assumption $w \in \mathbb{Z}_{\geq 0}$:
\begin{equation}
w = \frac{C}{R} \in \mathbb{Z}_{\geq 0}, \qquad R = \frac{C}{w} \tag{6.1}
\end{equation}
Dimensionless coordinates: $\tilde{x}^\mu = x^\mu / C \in \mathbb{Z}$.
\end{definition}

\subsection{Discrete Gnomonic Mapping}

\begin{definition}[Discrete gnomonic mapping]\label{def:discrete}
\begin{equation}
\tilde{L}^2 = 1 + w^2 \tilde{\sigma}^2, \qquad \tilde{\sigma}^2 = \sum_\mu \epsilon_\mu (\tilde{x}^\mu)^2 \tag{6.2}
\end{equation}
\begin{equation}
\tilde{Y}^\mu = \frac{\tilde{x}^\mu}{\tilde{L}}, \qquad \tilde{Y}^5 = \frac{1}{w\tilde{L}} \tag{6.3}
\end{equation}
\end{definition}

\subsection{Regularity}

\begin{theorem}[Signature-dependent regularity]\label{thm:reg}
\textbf{(i)} Euclidean: $\tilde{\sigma}^2 \geq 0 \Rightarrow \tilde{L}^2 \geq 1 > 0$. Globally regular.

\textbf{(ii)} Lorentzian: $\tilde{L}^2 = 0$ at
\begin{equation}
c^2t^2 = R^2 + x^2 + y^2 + z^2 \tag{6.6}
\end{equation}
This is the \textbf{cosmological horizon} in de Sitter spacetime~\cite{Guo2004}.
\end{theorem}

\begin{remark}[Diophantine equations on the horizon]
For $w = 1$ and isotropic points $\tilde{x} = \tilde{y} = \tilde{z} = s$, the horizon condition reduces to the Pell equation $\tilde{t}^2 - 3s^2 = 1$, with minimal solution $(s, \tilde{t}) = (1, 2)$.
\end{remark}

\begin{remark}[Discrete isotropy and dimensional constraints~{\rm\cite{Kihara2026b}}]\label{rem:fermat}
The author's previous work~\cite{Kihara2026b} determined the existence conditions for isotropic geodesics $l^2 + t^2 = Nx^2$ on $\mathbb{Z}^{1+N}$. By Fermat's two-square theorem: solutions exist if and only if every prime $p \equiv 3 \pmod{4}$ in $N$ appears with even exponent.

\begin{center}
\begin{tabular}{cccll}
\toprule
$N$ & $D=N\!+\!1$ & Prime factors & Solutions & Notes \\
\midrule
1 & 2 & $1$ & \checkmark & Pythagorean triples \\
2 & 3 & $2$ & \checkmark & \\
\textbf{3} & \textbf{4} & $3 \equiv 3$ (odd) & \textbf{None} & \textbf{Our universe} \\
4 & 5 & $2^2$ & \checkmark & Kaluza--Klein \\
5 & 6 & $5 \equiv 1$ & \checkmark & \\
9 & 10 & $3^2$ (even) & \checkmark & Superstrings \\
10 & 11 & $2 \times 5$ & \checkmark & M-theory \\
25 & 26 & $5^2$ & \checkmark & Bosonic strings \\
\bottomrule
\end{tabular}
\end{center}

$N = 3$ exhibits an arithmetic anomaly. $N = 4$ is the first extension admitting solutions, suggesting a fifth dimension is required for discrete isotropy.
\end{remark}

\begin{corollary}[Continuous limit]
As $C \to 0$: $\tilde{L}^2 \to l^2/R^2$ and $\tilde{Y}^\mu \to Rx^\mu/l = Y^\mu$, formally recovering~(1.1).
\end{corollary}

\subsection{Physical Identification of $C$}

\begin{hypothesis}[$C$ and Planck length]\label{hyp:planck}
\begin{equation}
C \sim \ell_P = \sqrt{\frac{\hbar G}{c^3}} \approx 1.616 \times 10^{-35}\,\text{m} \tag{6.5}
\end{equation}
\end{hypothesis}

%========================================
\section{Singularity Avoidance and Cosmological Horizons}
%========================================

\begin{theorem}[Euclidean regularity]
For fixed $R > 0$, $\Phi: \mathbb{R}^4 \to S^4(R)$ is $C^\infty$-regular on all of $\mathbb{R}^4$, since $l = \sqrt{R^2 + r^2} \geq R > 0$.
\end{theorem}

\begin{theorem}[Lorentzian regularity]
$\Phi$ is $C^\infty$-regular in the static patch $l^2 > 0$. The boundary $l^2 = 0$ defines the cosmological horizon.
\end{theorem}

\begin{proposition}[Singularity avoidance]
As $R \to 0$ ($r > 0$ fixed): $Y^5, Y^\mu \to 0$. The image degenerates but does not diverge.
\end{proposition}

\begin{theorem}[Bijectivity and information preservation]\label{thm:bij}
\textbf{(i)} Euclidean: $\Phi$ is a $C^\infty$-bijection onto the open hemisphere. Inverse $\Phi^{-1}$ exists.

\textbf{(ii)} Lorentzian: bijection within the static patch ($l^2 > 0$).
\end{theorem}

%========================================
\section{Consistency Checks}
%========================================

Under Hypothesis~\ref{hyp:planck} ($C = \ell_P$):

\textbf{Check 1:} $R_{\text{current}} = \sqrt{3/\Lambda_{\text{obs}}} \approx 1.65 \times 10^{26}\,\text{m} \approx 5.3\,\text{Gpc}$, consistent with the observable universe $R_H \approx 4.4 \times 10^{26}\,\text{m}$.

\textbf{Check 2:} $w = C/R_{\text{current}} \approx 10^{-61} \ll 1$. The present universe is near $w = 0$ (Minkowski).

\textbf{Check 3:} The gap between $w = 0$ ($\Lambda = 0$) and $w = 1$ ($\Lambda \sim 10^{122}\Lambda_{\text{obs}}$) is the discrete cosmological constant problem.

\begin{center}
\begin{tabular}{llll}
\toprule
Quantity & From $C = \ell_P$ & Known value & Status \\
\midrule
Planck length & $1.616 \times 10^{-35}$ m & $\ell_P$ & \checkmark\,(by def.) \\
Planck time & $5.39 \times 10^{-44}$ s & $t_P$ & \checkmark \\
$R_{\text{current}}$ & $1.65 \times 10^{26}$ m & $R_H \sim 4.4 \times 10^{26}$ m & \checkmark\,(same order) \\
\bottomrule
\end{tabular}
\end{center}

\begin{remark}[Avoiding circular reasoning]
The non-trivial check is that $R_{\text{current}} \sim R_H$ independently.
\end{remark}

%========================================
\section{Geometric Bijectivity and Information-Theoretic Implications}
%========================================

\begin{proposition}
Under Euclidean signature, $\Phi$ is a singularity-free $C^\infty$-bijection for all $R > 0$, preserving information in principle.
\end{proposition}

Under Lorentzian signature, information preservation is limited to the static patch. Application to the black hole information paradox requires quantum theory within the gnomonic framework (future work).

%========================================
\section{Outlook for Doctoral Research}
%========================================

\subsection{Extension to arbitrary dimensions}

The gnomonic framework extends to $S^n(R)$ for arbitrary $n$. The number-theoretic constraints (Remark~\ref{rem:fermat}) strongly constrain which dimensions admit isotropic discrete geodesics. The chain ``discrete isotropy $\to$ 5th dimension $\to$ Kaluza--Klein interpretation $\to$ emergence of electromagnetic force''~\cite{Kaluza1921,Klein1926} is the subject of a forthcoming paper.

\subsection{Fine-structure constant $\alpha$}
Connection to Wyler's formula~\cite{Wyler1971} via $SO(5,1) \to SO(4,2) \cong SU(2,2)/\mathbb{Z}_2$.

\subsection{Strong and weak forces}
Confinement ($V \sim \kappa r$) and flavor transitions require new extensions beyond this framework.

\subsection{Quantum theory, mass spectrum, and cosmology}
Discrete lattice structure connects to lattice field theory and loop quantum gravity. Derivation of $R(t)$ dynamics and Friedmann equation correspondence remain open.

%========================================
\section*{Summary of Results}
%========================================

\begin{center}
\begin{tabular}{lll}
\toprule
Item & Result & Status \\
\midrule
Lorentzian signature & Axioms A, B $\to$ $(-,+,+,+)$ & Axiom~3.1 \\
Quadratic metrics & Hurwitz theorem & Obs.~6.0 \\
Symmetry chain & $SO(5) \to SO(4,1) \to SO(3,1)$ & Prop.~\ref{prop:chain} \\
Angle upper bound & $540^\circ$ ($R$-independent) & Thm.~\ref{thm:angle} \\
Coupling constant & $w = C/R \in \mathbb{Z}_{\geq 0}$ & Def.~\ref{def:C} \\
Euclidean regularity & $\tilde{L}^2 \geq 1$ & Thm.~\ref{thm:reg}(i) \\
Lorentzian horizon & $l^2 = 0$ & Thm.~\ref{thm:reg}(ii) \\
Singularity avoidance & $C^\infty$ bijection & Thm.~\ref{thm:bij} \\
$R_{\text{current}} \sim R_H$ & Same order of magnitude & Consistency \\
\bottomrule
\end{tabular}
\end{center}

\begin{thebibliography}{99}

\bibitem{Paper1}
N.~Kihara, ``Geometric Formulation of 4-Dimensional Spacetime via Gnomonic Projection,'' 2026.
DOI: \href{https://doi.org/10.5281/zenodo.19427780}{10.5281/zenodo.19427780}.

\bibitem{doCarmo1992}
M.~P.~do Carmo, \textit{Riemannian Geometry}, Birkh\"{a}user, 1992, Chapter~4.

\bibitem{Wald1984}
R.~M.~Wald, \textit{General Relativity}, University of Chicago Press, 1984, Appendix~D.

\bibitem{GibbonsHawking1977}
G.~W.~Gibbons and S.~W.~Hawking, ``Action integrals and partition functions in quantum gravity,'' \textit{Phys.\ Rev.\ D}, \textbf{15}, 2752, 1977.

\bibitem{Guo2004}
H.-Y.~Guo, C.-G.~Huang, Z.~Xu, B.~Zhou, ``On Beltrami model of de Sitter spacetime,'' \textit{Mod.\ Phys.\ Lett.\ A}, \textbf{19}, 1701, 2004. arXiv:hep-th/0405137.

\bibitem{Kaluza1921}
T.~Kaluza, ``Zum Unit\"{a}tsproblem der Physik,'' \textit{Sitzungsber.\ Preuss.\ Akad.\ Wiss.}, 966--972, 1921.

\bibitem{Klein1926}
O.~Klein, ``Quantentheorie und f\"{u}nfdimensionale Relativit\"{a}tstheorie,'' \textit{Z.\ Phys.}, \textbf{37}, 895, 1926.

\bibitem{Wyler1971}
A.~Wyler, ``L'espace sym\'{e}trique du groupe des \'{e}quations de Maxwell,'' \textit{C.\ R.\ Acad.\ Sci.\ Paris}, \textbf{272}, 186, 1971.

\bibitem{MTW1973}
C.~W.~Misner, K.~S.~Thorne, J.~A.~Wheeler, \textit{Gravitation}, W.~H.~Freeman, 1973, \S13.5.

\bibitem{Spivak1999}
M.~Spivak, \textit{A Comprehensive Introduction to Differential Geometry}, Vol.~2, 3rd ed., Publish or Perish, 1999, Chapter~4.

\bibitem{Snyder1987}
J.~P.~Snyder, \textit{Map Projections: A Working Manual}, USGS Prof.\ Paper 1395, 1987, pp.~164--168.

\bibitem{Kihara2026b}
N.~Kihara, ``Solution of Positive Integer Solutions in Discrete Minkowski Intervals Using Gaussian Integers,'' GitHub: ai-chat-logs-open, 2026.

\bibitem{Hurwitz1898}
A.~Hurwitz, ``\"{U}ber die Composition der quadratischen Formen von beliebig vielen Variablen,'' \textit{Nachr.\ Ges.\ Wiss.\ G\"{o}ttingen}, 309--316, 1898.

\bibitem{Planck2020}
Planck Collaboration, ``Planck 2018 results.~VI.,'' \textit{Astron.\ Astrophys.}, \textbf{641}, A6, 2020.

\end{thebibliography}

\bigskip
\noindent\textit{This paper is confined to pure geometry. Physical claims and experimental predictions are left for future research.}\\
\textit{Last updated: 2026-04-06}

\end{document}
